\begin{titlepage}
\thispagestyle{empty}
\sffamily
\setlength{\parindent}{0pt}
\vspace*{0.19\textheight}

{\large\bfseries\color{DefinitionInk} Lean 形式化证明导读\par}
\vspace{1.2em}
{\color{DefinitionInk}\rule{27mm}{1pt}\par}
\vspace{1.7em}

{\fontsize{32pt}{45pt}\selectfont\bfseries\color{OxfordNavy}
 从 Theta 提升\par
 到广义 Riemann 猜想\par}
\vspace{0.8em}

{\Large\rmfamily\color{SlateInk}
 From Theta Lifts to the\par
 Generalized Riemann Hypothesis\par}
\vspace{2.6em}

{\normalsize\rmfamily
 本书根据 \href{https://github.com/lsernau/grh-lean}{\texttt{lsernau/grh-lean}}
 的可读版 Lean 证明及配套蓝图编写.\par}

\vfill
{\small\color{SlateInk}
 Lean 4.35.0\quad\textperiodcentered\quad
 Mathlib \texttt{56a00585}\par}
\end{titlepage}


\chapter*{内容提要}
\addcontentsline{toc}{chapter}{内容提要}
\markboth{内容提要}{内容提要}

本书以简化的 Lean 蓝图为基础，系统叙述 Riemann 猜想与 Dirichlet $L$ 函数广义 Riemann 猜想的完整论证. 该蓝图根据 Levent Alpöge 与 Claude 的 Lean~4 形式化证明整理而成. 正文从容许数据、加权素数幂余项与酉群上的阿代尔及代数 Theta 提升出发，建立四组核心算术与解析估计：短整源形式的阿基米德范数上界（定理 A(a)）、归一化提升族的局部格指数上界（定理 A(b)）、少量 Poincaré 平移对好目标形式空间的近乎张成性（定理 B），以及整形式提升链的余体积下界（定理 C）. 将这四组估计作用于同一个几乎张成目标空间的线性无关向量族，维数因子与公共主量精确抵消，从而导出主不等式 \eqref{eq:ch1-star}.

借助迭代重数 $m\asymp\sqrt{x}$ 带来的平方根量级节省，并结合固定模数的算术级数素数定理控制第二尺度上的误差，主不等式转化为惰性素数幂加权余项的单侧上界. 随后，利用二次扭转族的零点密度估计排除极点抵消，并经由 Landau 正性判据将单侧余项界转化为右半平面 $\RePart s>1/2$ 内的非零性. 最后，借助 Dirichlet $L$ 函数与 Riemann $\zeta$ 函数的函数方程，完成临界线定理及五条公开定理的证明.

\cleardoublepage
\pdfbookmark[0]{目录}{monograph-contents}
\tableofcontents
